% paper/sections/08b_pressure.tex
% §8.2「変密度射影法の段階的導出」
% ═══════════════════════════════════════════════════════════════════════════════
\subsection{変密度射影法の段階的導出}
\label{sec:projection_derivation}
% ═══════════════════════════════════════════════════════════════════════════════

第~\ref{sec:collocate}章（§~\ref{sec:helmholtz}）では等密度場を仮定した Helmholtz 分解で
射影法の幾何学的意味を説明した．
本節ではこれを\textbf{変密度場}へ拡張し，
なぜ可変係数演算子 $\bnabla\cdot(1/\rho\,\bnabla p)$ が現れるかを 3 段階で導出する．

本稿では van Kan（1986）の \textbf{IPC（Incremental Pressure Correction）法}~\cite{vanKan1986}を採用し，
予測段に前時刻圧力 $-\bnabla p^n$ を組み込む
（§~\ref{sec:time_int}，式~\eqref{eq:predictor_imex_bdf2}）．
これにより標準 Chorin 法~\cite{Chorin1968}の $\Ord{\Delta t}$ 分割誤差が $\Ord{\Delta t^2}$ に低減される
（二次精度射影法の系統解説は~\cite{BellColella1989,Guermond2006}）．
以下では圧力射影部分だけを抽出し，射影に用いる時間幅を
$\Delta t_{\mathrm{proj}}$ と書く
（IMEX-BDF2 では $\Delta t_{\mathrm{proj}}=\gamma\Delta t$，
BDF1 startup では $\Delta t_{\mathrm{proj}}=\Delta t$；§\ref{sec:ipc_derivation}）．

\subsubsection{第 1 段階：予測速度}

射影段で圧力項だけを分離し，対流・粘性・外力の時間離散済み寄与を
$\bm{R}_{\mathrm{proj}}^{n+1}$ とまとめると：
\begin{equation}
  \rho^{n+1}\frac{\bu^{n+1} - \bu^n}{\Delta t_{\mathrm{proj}}}
  = \bm{R}_{\mathrm{proj}}^{n+1}
  - \bnabla p^{n+1}
  \label{eq:NS_full}
\end{equation}
IPC 法では前時刻圧力 $-\bnabla p^n$ を陽的に加えて\textbf{予測速度} $\bu^*$ を定義する：
\begin{equation}
  \rho^{n+1}\frac{\bu^* - \bu^n}{\Delta t_{\mathrm{proj}}}
  = \bm{R}_{\mathrm{proj}}^{n+1} - \bnabla p^n
  \quad\Longrightarrow\quad
  \bu^* = \bu^n
  + \frac{\Delta t_{\mathrm{proj}}}{\rho^{n+1}}
    \bigl(\bm{R}_{\mathrm{proj}}^{n+1} - \bnabla p^n\bigr)
  \label{eq:predictor}
\end{equation}
$\bu^*$ は一般に発散フリーではない（$\bnabla\cdot\bu^* \neq 0$）．

\subsubsection{第 2 段階：圧力補正}

式~\eqref{eq:NS_full} から式~\eqref{eq:predictor} を引くと，
圧力増分 $\delta p \equiv p^{n+1} - p^n$ と速度補正の関係が得られる：
\begin{equation}
  \rho^{n+1}\frac{\bu^{n+1} - \bu^*}{\Delta t_{\mathrm{proj}}} = -\bnabla(\delta p)
  \quad\Longrightarrow\quad
  \bu^{n+1} = \bu^* - \frac{\Delta t_{\mathrm{proj}}}{\rho^{n+1}}\bnabla(\delta p)
  \label{eq:corrector}
\end{equation}

\subsubsection{第 3 段階：発散フリー条件から PPE を導く}

$\bu^{n+1}$ に非圧縮条件 $\bnabla\cdot\bu^{n+1}=0$ を課す．
式~\eqref{eq:corrector} の両辺に発散を作用させると：
\begin{equation}
  \bnabla\cdot\!\left(\frac{1}{\rho^{n+1}}\bnabla(\delta p)\right)
  = \frac{1}{\Delta t_{\mathrm{proj}}}\,\bnabla\cdot\bu^*
  \label{eq:PPE_varrho}
\end{equation}
これが変密度 PPE（IPC 増分形式）である．
一様密度（$\rho=\mathrm{const}$）では
$\bnabla^2(\delta p) = \rho\,\bnabla\cdot\bu^*/\Delta t_{\mathrm{proj}}$ に帰着する．

\noindent\textbf{補足：}
外力 $\bm{F}_\text{ext}^{n+1}$ は重力項を含む．
表面張力の扱いは経路で分かれる：一括 smoothed-Heaviside PPE の縮約検証では
CSF 体積力 $(\kappa_{lg}^{n+1}\bnabla\psi^{n+1})/We$ を予測段に保持するが，
本稿の標準圧力ジャンプ経路では CSF 体積力としては入れず，
$j_{gl}=p_g-p_l=-\sigma\kappa_{lg}$ から作る毛管補正量を
第 3 段階の PPE/補正段に組み込む
（§~\ref{sec:split_ppe_pressure_jump_formulation}）．
密度 $\rho^{n+1}$ は，界面輸送により新時刻の配置として与えられている．
本節の導出は Backward Euler 形式で統一したが，
対流項は IMEX-BDF2（式~\eqref{eq:predictor_imex_bdf2}），
粘性項は implicit-BDF2（$\Ord{\Delta t^2}$，§~\ref{sec:time_viscous}）を用いる．
高密度比の分相 PPE 解法における圧力勾配の評価には
HFE による場延長を要する（§~\ref{sec:field_extension}，§~\ref{sec:split_ppe_recovery}）．

\noindent\textbf{PPE の離散化方針：}
以下の演算子形式は圧力増分 $\delta p$ および圧力 $p$ のいずれにも適用可能であるため，
一般性のため $p$ で表記する．
一括 PPE 縮約で式~\eqref{eq:PPE_varrho} の可変係数演算子を product-rule で展開すると：
\begin{equation}
  \bnabla\cdot\!\left(\frac{1}{\rho}\bnabla p\right)
  = \frac{1}{\rho}\bnabla^2 p - \frac{\bnabla\rho}{\rho^2}\cdot\bnabla p
  \label{eq:varrho_product_rule}
\end{equation}
この形式は高次 CCD 系の演算子で表せる（§~\ref{sec:ccd_poisson_matrix}）．
ただし高密度比・低 We の標準経路では，product-rule 展開した一括場ではなく，
分相圧力ジャンプ PPE と欠陥補正法（§~\ref{sec:defect_correction_ppe}）で
同じ有効面空間 $D_fG_A$ の固定点を定める．
有限体積法との等価性・調和平均の根拠は §~\ref{sec:fvm_face_coeff} 参照．
仮想時間刻みとゲージ固定の補足は付録~\ref{sec:ppe_pseudotime} 参照．

ただし，式~\eqref{eq:PPE_varrho} をどの線形系として解くかだけでは，
静止界面における力の釣り合いは保証されない．
PPE 内の圧力勾配，速度補正に用いる圧力勾配，および界面力が
同じ面評価位置と係数を共有することが，後続の Balanced--Force 条件の主題である．

\subsubsection{圧力反力としての面勾配：PPE だけでは代表が決まらない}
\label{sec:variational_pressure_reaction}

圧力射影を面速度空間で書くと，補正段は
\begin{equation}
  \bu_f^{n+1}
  =
  \bu_f^\ast
  -
  \Delta t_{\mathrm{proj}}\,
  a_f(p;c),
  \qquad
  a_f(p;c)=G_Ap-c_f
  \label{eq:subtractive_face_pressure_reaction}
\end{equation}
である．ここで $c_f$ は既知の毛管補正量または界面ジャンプ補正量であり，
$G_Ap$ は圧力が面速度へ及ぼす反力である．
この $G_A$ は，単に $D_fG_Ap$ が PPE 左辺を作るという条件だけでは決まらない．
圧力が非圧縮制約のラグランジュ乗数であるためには，面運動量計量 $M_A$ と
圧力節点計量 $W_p$ に対して
\begin{equation}
  \langle G_Ap,w_f\rangle_{M_A}
  +
  \langle p,D_fw_f\rangle_{W_p}
  =0
  \qquad
  \forall p,\ \forall w_f
  \label{eq:pressure_reaction_green_identity}
\end{equation}
を満たさなければならない（符号は式~\eqref{eq:subtractive_face_pressure_reaction}
の減算型補正に合わせたもの）．
したがって
\begin{equation}
  G_A=-M_A^{-1}D_f^TW_p,
  \qquad
  L_Ap\equiv D_fG_Ap
  \label{eq:variational_pressure_operator_pair}
\end{equation}
が，固定した $D_f,M_A,W_p$ に対する圧力反力とスカラー PPE の組である．
この組を用いると，圧力補正による仕事は
\begin{equation}
  \langle G_Ap,G_Ap\rangle_{M_A}
  =
  -\langle p,L_Ap\rangle_{W_p}
  \label{eq:pressure_reaction_energy_identity}
\end{equation}
として同じ離散内積で読める（$L_A$ の符号は
式~\eqref{eq:variational_pressure_operator_pair} の発散向きに従う）．

ここで重要なのは，$D_fG_Ap$ が同じでも $G_Ap$ の代表は一意ではない点である．
もし
\[
  \widetilde Gp = G_Ap+Zp,\qquad D_fZp=0,\qquad Zp\ne0
\]
なら，$\widetilde G$ は同じ PPE 右辺を満たし得るが，
式~\eqref{eq:pressure_reaction_green_identity} の圧力仕事を一般には満たさない．
この発散ゼロの余剰成分は，静的液滴では寄生流れとして，
非平衡界面では毛管駆動の過小または過大評価として現れる．
したがって本稿の圧力ジャンプ標準経路では，
補正段の圧力面加速度を
式~\eqref{eq:pressure_reaction_green_identity} の変分随伴代表に限定し，
高次コンパクト勾配のような発散同値な代表は，
この随伴性が証明された場合を除き，圧力反力としては用いない．
保存形共通フラックス経路では，この圧力射影の $D_f$ が
界面・密度・運動量輸送の $D_f$ でもなければならない．
射影が $D_f\bu_f=0$ を満たすという主張は，その同じ $D_f$ で
$q,m,\bm{p}_m$ が更新されるときにだけ質量輸送の非圧縮性を意味する．
したがって圧力反力の随伴性と共通フラックス台帳は，
どちらも同じ面速度空間を固定するための条件である．

\subsubsection{非一様格子での面平均補正}
\label{sec:fvm_ccd_corrector}

均一格子・周期境界条件では，PPE 内の勾配と速度補正の勾配に同じ
CCD 演算子を用いれば P-2（PPE 側 $G_h$ と補正側 $G_h$ の共有）が成立する．
一方，界面適合非一様格子と壁面境界条件では，PPE が面フラックス
$D_h(A_f G_h p)$ を有限体積法的に積分するのに対し，節点 CCD の速度補正は
節点メトリックで勾配を評価するため，面メトリックとの不一致が残る．

この経路では速度補正を面平均演算子
$\mathcal{G}^\text{adj}$ に切り替え，PPE の離散拡散項も
同じ随伴ペア $D_h^\text{bf}=-(\mathcal{G}^\text{adj})^\ast$ から組み立てる．
これにより PPE と補正が同じ面フラックス空間を共有し，
非一様格子での有限体積法--CCD メトリック不整合を避ける．
圧力ジャンプ条件を併用する場合も同じ面フラックス空間で
$G_{\Gamma}^{\mathrm{adj}}p=G^{\mathrm{adj}}p-B^{\mathrm{adj}}(j_{gl})$ とし，
$B^{\mathrm{adj}}(j_{gl})$ を面平均勾配と同じ局所物理距離 $H_f$ で定義する．
この面平均条件は本節で定義し，静止液滴での安定化効果は
§~\ref{sec:bf_static_droplet} で検証する．

\phantomsection
\label{sec:dccd_pressure_nofilt}% backward-compat: 外部参照保持
